\clearpage
\SourceLeaf{020b}
\begin{center}
\pdfbookmark[2]{2.}{article-002}2.
\end{center}
Leges modo expositae ab iis, quas \textsc{Kepler}us noster detexit, aliter non differunt, nisi quod in forma ad omnia sectionum conicarum genera patente exhibitae sunt, actionisque corporis moti in Solem, a qua pendet factor $\BareRoot{(1+\mu)}$, ratio est habita. Si has leges tamquam phaenomena ex innumeris atque indubiis observationibus depromta consideramus, geometria docebit, qualis actio
\SourceLeaf{021}in corpora circa Solem mota ab hoc exerceri debeat, ut ista phaenomena perpetuo producantur. Hoc modo invenitur, Solis actionem in corpora ambientia perinde se exercere, ac si vis attractiva, cuius intensitas quadrato distantiae reciproce proportionalis esset, corpora versus centrum Solis propelleret. Quodsi vero vice versa a suppositione talis vis attractivae tamquam principio proficiscimur, phaenomena illa ut consequentiae necessariae inde derivantur. Hic leges tantum enarravisse sufficiat, quarum nexui cum principio gravitationis hoc loco eo minus opus erit immorari, quum post summum \textsc{Newton} auctores plures hoc argumentum tractaverint, interque eos ill. \textsc{Laplace} in opere perfectissimo, Mécanique Céleste, tali modo, ut nihil amplius desiderandum reliquerit.

\begin{center}
\pdfbookmark[2]{3.}{article-003}3.
\end{center}

Disquisitiones circa motus corporum coelestium, quatenus fiunt in sectionibus conicis, theoriam completam huius curvarum generis neutiquam postulant: quin adeo unica aequatio generalis nobis sufficiet, cui omnia superstruantur. Et quidem maxime e re esse videtur, eam ipsam eligere, ad quam tamquam aequationem characteristicam deferimur, dum curvam secundum attractionis legem descriptam investigamus. Determinando scilicet quemvis corporis locum in orbita sua per distantias $x, y$ a duabus rectis in plano orbitae ductis atque in centro Solis i.\,e. in altero curvae foco sub angulis rectis se secantibus, et denotando insuper corporis distantiam a Sole (positive semper accipiendam) per $r$, habebimus inter $r, x, y$ aequationem linearem $r+\alpha x+\beta y=\gamma$, in qua $\alpha, \beta, \gamma$ quantitates constantes exprimunt, et quidem $\gamma$ quantitatem natura sua semper positivam. Mutando rectarum, ad quas distantiae $x, y$ referuntur, situm per se arbitrarium, si modo sub angulis rectis se intersecare perseverent, manifesto forma aequationis valorque ipsius $\gamma$ non mutabuntur, $\alpha$ et $\beta$ autem alios aliosque valores nanciscentur, patetque, situm illum ita determinari posse, ut $\beta$ evadat $=0$, $\alpha$ autem saltem non negativa. Hoc modo scribendo pro $\alpha, \gamma$ resp. $e, p$, aequatio nostra induit formam $r+ex=p$. Recta, ad quam tunc distantiae $y$ referuntur, \emph{linea apsidum} vocatur, $p$ \emph{semiparameter}, $e$ \emph{excentricitas}; sectio conica denique \emph{ellipsis}, \emph{parabolae} vel \emph{hyperbolae} nomine distinguitur, prout $e$ unitate minor, unitati aequalis, vel unitate maior est.

\SourceLeaf{022}Ceterum facile intelligitur, situm lineae apsidum per conditiones traditas plene determinatum esse, unico casu excepto, ubi tum $\alpha$ tum $\beta$ iam per se erant $=0$; in hoc casu semper fit $r=p$, ad quascunque rectas distantiae $x, y$ referantur. Quoniam itaque habetur $e=0$, curva (quae erit circulus) secundum definitionem nostram ellipsium generi annumeranda est, id vero singulare habet, quod apsidum positio prorsus arbitraria manet, siquidem istam notionem ad hunc quoque casum extendere placet.
